%% Question bank chapter file - generated by bank_tools.py build (do not edit while the bank is still being built:
%% it is regenerated from the per-test data in Question Bank\data). Every question is one bankq block:
%%   \begin{bankq}[id=..., source={...}, paper=.., qno=.., type=..., subtopic={...}, answer={...}, answer_src=key|solved,
%%                 working={...}, difficulty=1-10, diff_parts={K S I C F}, diff_note={...},
%%                 group=<passage id or empty>, also_in={duplicate copies in other tests}]
%% difficulty = min(10, K+S+I+C+F-1): K concepts 1-3, S steps 1-4, I insight 0-2, C computation 0-2, F format 0-1
%%   ...question LaTeX, options, figures (\bfig / \bfigi -> img/)...
%%   \end{bankq}
%% Passages shared by several questions are bankpassage blocks; their questions carry group=<passage id>.
\documentclass[10pt]{article}
\usepackage[a4paper,margin=1.8cm]{geometry}
\usepackage[T1]{fontenc}
\usepackage{amsmath,amssymb,mathtools}
\usepackage{graphicx}
\usepackage[export]{adjustbox}
\usepackage{xcolor,needspace,keyval}
\usepackage{enumitem}
\usepackage{booktabs,array,tabularx}
\usepackage[hidelinks]{hyperref}
\setlength{\parindent}{0pt}
\setlength{\parskip}{3pt}
\makeatletter
\define@key{bq}{id}{\def\bq@id{#1}}\define@key{bq}{source}{\def\bq@src{#1}}\define@key{bq}{paper}{\def\bq@paper{#1}}
\define@key{bq}{qno}{\def\bq@qno{#1}}\define@key{bq}{type}{\def\bq@type{#1}}\define@key{bq}{subtopic}{\def\bq@sub{#1}}
\define@key{bq}{answer}{\def\bq@ans{#1}}\define@key{bq}{answer_src}{\def\bq@asrc{#1}}\define@key{bq}{working}{\def\bq@work{#1}}
\define@key{bq}{group}{\def\bq@group{#1}}\define@key{bq}{also_in}{\def\bq@also{#1}}
\define@key{bq}{difficulty}{\def\bq@diff{#1}}\define@key{bq}{diff_parts}{\def\bq@dparts{#1}}\define@key{bq}{diff_note}{\def\bq@dnote{#1}}
\newenvironment{bankq}[1][]{\def\bq@work{}\def\bq@also{}\def\bq@group{}\def\bq@diff{?}\def\bq@dparts{}\def\bq@dnote{}\setkeys{bq}{#1}%
  \par\needspace{4\baselineskip}\medskip\noindent{\small\textbf{\bq@id}\hfill\color{gray}[\bq@src, Paper \bq@paper, Q\bq@qno]
  \bq@type\ $\cdot$ \bq@sub\ $\cdot$ \textbf{D\bq@diff}}\par\nopagebreak\relax}
  {\par{\small\color{teal}Answer: \textbf{\bq@ans}\ifx\bq@work\@empty\else\ (solved: \bq@work)\fi
  \ifx\bq@also\@empty\else\par\color{gray}Also in: \bq@also\fi}\par}
\define@key{bp}{id}{\def\bp@id{#1}}\define@key{bp}{source}{\def\bp@src{#1}}\define@key{bp}{questions}{\def\bp@qs{#1}}
\newenvironment{bankpassage}[1][]{\setkeys{bp}{#1}\par\needspace{4\baselineskip}\bigskip\noindent
  {\small\textbf{Passage \bp@id}\hfill\color{gray}[\bp@src; questions \bp@qs]}\par\nopagebreak\relax\itshape}{\par}
\makeatother
\newcommand{\bfig}[2]{\par\smallskip{\centering\includegraphics[width=#1pt,max width=0.62\linewidth,max height=0.32\textheight]{img/#2}\par}\smallskip}
\newcommand{\bfigi}[2]{\includegraphics[width=#1pt,max width=\linewidth,max height=0.32\textheight,valign=c]{img/#2}}
\newcommand{\bmissing}[1]{\textit{[figure missing in the source conversion: \detokenize{#1}]}}
\newenvironment{qoptions}{\begin{itemize}[leftmargin=1.8em,itemsep=1pt,topsep=2pt,label={}]}{\end{itemize}}
\newcommand{\opt}[2]{\item[(#1)] #2}
\providecommand{\fwstart}{}\providecommand{\fwend}{}\providecommand{\qfullsplit}{}
\newcolumntype{L}[1]{>{\raggedright\arraybackslash}p{#1}}
\newcolumntype{C}[1]{>{\centering\arraybackslash}p{#1}}
\newcolumntype{R}[1]{>{\raggedleft\arraybackslash}p{#1}}
\providecommand{\sectionhdr}[1]{}\providecommand{\subsectionhdr}[1]{}


\begin{document}
{\Large\bfseries Statistics --- Question Bank}\par\smallskip
{\small Mathematics chapter M09. Questions from all converted test papers, grouped by source series. Answers: printed key where it matches the paper, otherwise solved (answer\_src=solved, with working). Exact repeats across papers appear once, with also\_in listing the other copies.}\par\medskip\hrule

\section*{Score Advanced papers}
\begin{bankq}[id=SA05-P1-MATH-S2-Q09, source={Score Advanced - Test 05 (15 Jul 2021)}, paper=1, qno=9, type=numerical, subtopic={Variance and standard deviation}, answer={8.665}, answer_src=key, difficulty=5, diff_parts={K2 S1 I1 C1 F1}, diff_note={concepts: corrected variance, variance and standard deviation | steps: replace 3 by 4 in the sums}]
Let $\displaystyle\sum_{i=1}^{9}x_i^2 = 800$ and $\displaystyle\sum_{i=1}^{9}x_i = 80$ where x$_i$ is i$^{\text{th}}$ observation. If one observation that was 3, was found to be wrong and replaced by corrected value 4 then corrected variance is
\end{bankq}

\section*{Nurture SRG papers}
\begin{bankq}[id=NSRG02-P1-MATH-S1-Q02, source={Nurture SRG - Test 02 (06 Dec 2023)}, paper=1, qno=2, type=single, subtopic={Variance and standard deviation}, answer={C}, answer_src=solved, working={For binomial frequencies the variable is $\mathrm{Bin}(n,\tfrac12)$ with variance $\dfrac n4$}, difficulty=5, diff_parts={K1 S2 I2 C1 F0}, diff_note={concepts: variance of a binomial frequency distribution | steps: mean from binomial sums; mean of squares; variance}]
If the variable takes values $0,1,2,3,\dots,n$ with frequencies ${}^nC_0,{}^nC_1,{}^nC_2,\dots,{}^nC_n$ respectively, then variance is
\begin{qoptions}
\opt{A}{$\dfrac n3$}
\opt{B}{$\dfrac{2n}5$}
\opt{C}{$\dfrac n4$}
\opt{D}{None of these}
\end{qoptions}
\end{bankq}
\begin{bankq}[id=NSRG02-P1-MATH-S1-Q03, source={Nurture SRG - Test 02 (06 Dec 2023)}, paper=1, qno=3, type=single, subtopic={Variance and standard deviation}, answer={B}, answer_src=solved, working={$\sum x^2=5(18+150^2)=112590$; new sum of squares $136926$, new mean 151, variance $22821-22801=20$}, difficulty=5, diff_parts={K1 S2 I1 C2 F0}, diff_note={concepts: variance after adding an observation | steps: sum and sum of squares; new sums; new variance}]
Average height and variance of 5 students in a class is 150cm and 18cm$^2$ respectively. A new student whose height is 156 cm is added to the group. The new variance in cm$^2$ is
\begin{qoptions}
\opt{A}{16}
\opt{B}{20}
\opt{C}{22}
\opt{D}{14}
\end{qoptions}
\end{bankq}
\begin{bankq}[id=NSRG02-P1-MATH-S2-Q07, source={Nurture SRG - Test 02 (06 Dec 2023)}, paper=1, qno=7, type=multi, subtopic={Mean, median, mode}, answer={A,D}, answer_src=solved, working={Sorted numbers $x-y,y,x,2x+y$; median $\tfrac{x+y}2=10$, range $x+2y=28$ give $x=12,y=8$, mean 14}, difficulty=7, diff_parts={K1 S3 I2 C1 F1}, diff_note={concepts: mean median and range of four expressions | steps: order the numbers; median equation; range equation; solve; mean}]
If the median and range of four numbers $\{x,\,y,\,2x+y,\,x-y\}$, where $0<y<x<2y$, are 10 and 28 respectively, then which of the following is/are CORRECT?
\begin{qoptions}
\opt{A}{$x+y=20$}
\opt{B}{$x-y=6$}
\opt{C}{Mean of the numbers is 18}
\opt{D}{Mean of the numbers is 14}
\end{qoptions}
\end{bankq}
\begin{bankq}[id=NSRG02-P1-MATH-S3-Q05, source={Nurture SRG - Test 02 (06 Dec 2023)}, paper=1, qno=5, type=numerical, subtopic={Variance and standard deviation}, answer={3}, answer_src=solved, working={$\sigma^2+(\mu+2)^2=18$, $\sigma^2+(\mu-2)^2=10$ give $\mu=1$, $\sigma^2=9$, $\sigma=3$}, difficulty=6, diff_parts={K1 S2 I2 C1 F1}, diff_note={concepts: mean and standard deviation from variances about two points | steps: variance about a point; subtract equations; mean; standard deviation}]
The variance about $-2$ and 2 are 18 and 10 respectively. Find standard deviation of the set of observation
\end{bankq}
\begin{bankq}[id=NSRG03-P1-MATH-S3-Q06, source={Nurture SRG - Test 03 (04 Feb 2024)}, paper=1, qno=6, type=numerical, subtopic={Variance and standard deviation}, answer={1.25}, answer_src=solved, working={$\sum x=100$, mean 5, $\sigma^2=5$, so $\lambda=5$ and $\lambda/4=1.25$}, difficulty=6, diff_parts={K1 S2 I2 C1 F1}, diff_note={concepts: variance after shifting observations | steps: mean from the shifted data; variance from the second condition; scale}]
There are 20 observations $x_1,x_2,\dots,x_{20}$. If mean of $(x_i+i)\ \forall\ i=1,2,3,\dots20$ is 15.5 and $\displaystyle\sum_{i=1}^{20}(x_i-5)^2=100$. If variance of $(x_i-2)\forall\ i=1,2,3\dots,20$ is $\lambda$, then $\dfrac\lambda4$ is equal to
\end{bankq}
\begin{bankq}[id=NSRG03-P2-MATH-S3-Q06, source={Nurture SRG - Test 03 (04 Feb 2024)}, paper=2, qno=6, type=numerical, subtopic={Variance and standard deviation}, answer={3}, answer_src=solved, working={$\sigma^2+(\mu+2)^2=18$ and $\sigma^2+(\mu-2)^2=10$ give $\mu=1$, $\sigma=3$}, difficulty=6, diff_parts={K1 S2 I2 C1 F1}, diff_note={concepts: standard deviation from deviations about two points | steps: mean square deviation about a point; subtract; mean; variance}]
The mean square deviation of a set of $n$ observations $x_1,x_2,\dots,x_n$ about a point $c$ is defined as $\dfrac1n\displaystyle\sum_{i=1}^n\left(x_i-c\right)^2$. The mean square deviation about $-2$ and 2 are 18 and 10 respectively, then standard deviation of this set of observations is :
\end{bankq}

\section*{Star Review (Other) papers}
\begin{bankq}[id=STAR03-P1-MATH-S2-Q05, source={Star Review (Other) - Test 03 (04 Dec 2023)}, paper=1, qno=5, type=multi, subtopic={Variance and standard deviation}, answer={C}, answer_src=key, difficulty=5, diff_parts={K1 S2 I1 C1 F1}, diff_note={concepts: variance from shifted sums | steps: mean from the shifted sum; shifted sum of squares; variance formula}]
If $\displaystyle\sum_{i=1}^{10}(x_i-5)=5$ and $\displaystyle\sum_{i=1}^{10}(x_i-5)^2=125$, variance of $x_1,x_2,\ldots\ldots x_{10}$ will be-
\begin{qoptions}
\opt{A}{$\dfrac94$}
\opt{B}{$\dfrac{25}{4}$}
\opt{C}{$\dfrac{49}{4}$}
\opt{D}{$\dfrac{81}{4}$}
\end{qoptions}
\end{bankq}
\end{document}
